\paragraph{What is new.} Earlier CFL methods cluster whatever the clients send (updates, losses, data statistics) with a similarity whose scale depends on sample size, on training progress and on noise. The number of clusters then depends on a hyperparameter that has no fixed meaning. DisCo-CFL defines the clustering target precisely (class-conditional concept equivalence) and estimates a similarity whose population value is known in advance for clients that share a concept, namely $1$. The invariance, the automatic choice of $K$, the size-independent recovery, the robustness to DP noise and the explanation guarantees are all consequences of this calibration.

\paragraph{Limitations.}
(i) The class support $S_i$ of each client and the per-class half sizes are revealed to the server; the DP guarantee protects individual records given the labels, not the label histogram.
(ii) The analysis assumes independent sub-Gaussian per-sample signatures and a fixed reference model. It does not cover the attachment heuristics, and the constants are not optimised. The upper bound exceeds the lower bound by a factor $1/\Delta$.
(iii) Concept shifts that leave the reference-model gradient unchanged are invisible by construction (Assumption~\ref{ass:sep}). An example is a $180^\circ$ rotation of nearly symmetric Fashion-MNIST items; for classification such shifts are also largely harmless.
(iv) The signatures are only as informative as the reference model. On CIFAR-10 a 10-round warm-up gives a weak model and weak separation, and 30 rounds are needed for good clustering. An adaptive warm-up that stops when the split-half reliabilities and the between-cluster gaps stabilise is a natural extension.
(v) Clients that hold none of the discriminative classes are compatible with several groups. DisCo-CFL may place them in any of these groups, which is harmless for training but lowers ARI.
(vi) Clustering is one-shot at the reference model. Concept drift after clustering would require recomputing signatures, although the newcomer procedure already detects clients whose concept is unseen.
(vii) Our main experiments use small CNNs and synthetic concept groups built from public image datasets, as is standard in the CFL literature~\citep{benali2026survey}. With a pretrained ResNet-18 (Section~\ref{sec:real}) the method works on CIFAR-100 with label swaps and, partly, on PACS, but it does not separate the groups on Office-Home or on rotated CIFAR-100. We tried a signature from the whole last block and a longer warm-up on one seed, without success, so the cause is not settled; the identifiability condition (Assumption~\ref{ass:sep}) is the likely one, but we have not verified it. In both failures the oracle gains less than one point over FedAvg, so the cost in accuracy is small, but a setting where clustering matters and DisCo-CFL misses it is not excluded. The realistic benchmarks have four groups and 40 clients, trained from one pretrained network; larger federations and other architectures remain to be tested.
(viii) On natural writer heterogeneity (FEMNIST) the method forms many small clusters (47 for 100 writers). It reaches the highest mean accuracy there, but the cluster models generalise less to other writers than the global model, and we have no way to tell whether the splits correspond to real concept differences.
